\chapter{Matriz de consistencia}
\label{anexo:c}

La matriz de consistencia vincula \textbf{problema tecnológico}, \textbf{objetivo general
(técnico)}, las \textbf{cinco fases} del modelo S/E/C, variables, hipótesis, indicadores y
controles. Fue ajustada según retroalimentación docente PI-10 (eliminación de piloto Payment
Execution; Fase~1 = diseño del procedimiento, no experimentación) y PI-11 (separación de niveles
del problema).

\section{Metadatos}

\begin{description}
  \item[Situación real (problemática):] Síntomas del sector bancario: fragmentación OpenAPI,
        costos de integración, casos documentados de desalineación entre sistemas (Cap.~1 §1.2).
        \textbf{No} es el problema de las filas de la matriz.
  \item[Situación subyacente (abstracto):] Correspondencia entre estructuras OpenAPI y
        \ServiceDomain{} BIAN; objeto y unidad de análisis (Cap.~1 §1.3--§1.4).
  \item[Problema tecnológico (matriz):] \textbf{No existe} un método reproducible que cuantifique
        en qué medida un contrato OpenAPI bancario se alinea con un \ServiceDomain{} BIAN y bajo
        qué tipologías se desalinea. Pregunta ancla: ¿en qué medida… y cómo cuantificarlo mediante
        el modelo S/E/C $\rightarrow$ \AlignmentScore{}? (Cap.~1 §1.3; distinto de la brecha del
        estado del arte, Cap.~3 §3.6).
  \item[Objetivo superior:] Consecuencia en el mundo real si el artefacto operara: gobernanza y
        reducción de retrabajo en integraciones (Cap.~2 §2.2).
  \item[Objetivo general (técnico):] Diseñar el modelo de alineación y el procedimiento que
        calculen \SemanticScore{} (S), \StructuralScore{} (E), \CoverageScore{} (C) y
        \AlignmentScore{} (Cap.~2 §2.3).
\end{description}

\section{Resumen por fase}

\small
\begin{longtable}{@{}p{0.14\textwidth}p{0.26\textwidth}p{0.26\textwidth}p{0.28\textwidth}@{}}
\caption{Matriz de consistencia — resumen por fase (Anexo C)} \label{tab:anexo-c} \\
\toprule
\textbf{Fase} & \textbf{Objetivo} & \textbf{Producto verificable} & \textbf{Hipótesis (síntesis)} \\
\midrule
\endfirsthead
\multicolumn{4}{l}{\textit{(continúa)}} \\
\midrule
\textbf{Fase} & \textbf{Objetivo} & \textbf{Producto verificable} & \textbf{Hipótesis (síntesis)} \\
\midrule
\endhead
\midrule
\multicolumn{4}{r}{\textit{(continúa en la siguiente página)}} \\
\endfoot
\bottomrule
\endlastfoot
Fase 1 & Diseñar el modelo de alineación OpenAPI $\leftrightarrow$ BIAN, definiendo entradas, componentes de score (S, E, C), fórmula de agregación y criterios de clasificación. & Modelo metodológico documentado: adquisición → normalización → SemanticScore (S) + StructuralScore (E) + CoverageScore (C) → AlignmentScore. & Un modelo estructurado por componentes S, E y C permite calcular de forma reproducible un score de alineación interpretable entre OpenAPI y BIAN. \\
Fase 2 & Transformar contratos OpenAPI y artefactos BIAN en representaciones normalizadas y comparables sobre un modelo intermedio común. & Conjunto de representaciones normalizadas versionadas (paths, operations, schemas, behaviors, business objects) listas para el cálculo de scores. & La normalización de OpenAPI bancario y BIAN en un modelo intermedio común habilita el cálculo sistemático de los scores S, E y C. \\
Fase 3 & Calcular el StructuralScore (E): correspondencia estructural entre componentes OpenAPI y artefactos BIAN normalizados. & Matriz de correspondencias estructurales y StructuralScore (E) normalizado [0, 1]. & Las técnicas de schema matching producen un StructuralScore que refleja la correspondencia entre la organización REST y el modelo BIAN. \\
Fase 4 & Calcular el SemanticScore (S): similitud semántica entre conceptos de negocio BIAN y componentes del contrato OpenAPI. & Matriz de similitud semántica y SemanticScore (S) normalizado [0, 1]. & Las métricas de similitud semántica producen un SemanticScore que discrimina conceptos alineados de desalineaciones semánticas o de modelado. \\
Fase 5 & Calcular CoverageScore (C), obtener AlignmentScore = $\alpha$--S + $\beta$--E + $\gamma$--C ($\alpha$ + $\beta$ + $\gamma$ = 1) y clasificar el nivel de alineación del contrato OpenAPI. & CoverageScore (C); AlignmentScore; clasificación (Alta $\geq$0,80 -- Media 0,60–0,79 -- Baja 0,40–0,59 -- Nula <0,40); informe de tipologías de desalineación. & La agregación ponderada de S, E y C produce un AlignmentScore interpretable y clasificable por niveles de alineación del contrato OpenAPI bancario. \\
\end{longtable}

\normalsize

\section{Trazabilidad Cap.~2}

\begin{table}[H]
  \centering
  \caption{Vinculación fases del modelo $\rightarrow$ objetivos específicos}
  \label{tab:anexo-c-trazabilidad}
  \small
  \begin{tabular}{@{}lll@{}}
    \toprule
    \textbf{Fase modelo} & \textbf{OE} & \textbf{Score / producto} \\
    \midrule
    Fase 1 — Diseño & Marco OG & Modelo S/E/C documentado \\
    Fase 2 — Normalización & OE1 & Representaciones comparables \\
    Fase 3 — Matching estructural & OE2 & \StructuralScore{} (E) \\
    Fase 4 — Similitud semántica & OE3 & \SemanticScore{} (S) \\
    Fase 5 — Integración & OE4 & \CoverageScore{} (C) + \AlignmentScore{} \\
    \bottomrule
  \end{tabular}
\end{table}
